\hypertarget{fonctions}{%
\section{Fonctions}\label{fonctions}}

\hypertarget{appels-de-fonctions}{%
\subsection{Appels de fonctions}\label{appels-de-fonctions}}

Dans le contexte de la programmation, une fonction est une séquence
d\textquotesingle instructions nommée qui effectue un calcul. Lorsque
vous définissez une fonction, vous spécifiez son nom et la séquence
d\textquotesingle instructions. Plus tard, vous pouvez « appeler » la
fonction par son nom.

Nous avons déjà vu un exemple d\textquotesingle appel de fonction :

\begin{Shaded}
\begin{Highlighting}[]
\OperatorTok{\textgreater{}\textgreater{}\textgreater{}} \BuiltInTok{print}\NormalTok{(}\DecValTok{32}\NormalTok{)}
\DecValTok{32}
\end{Highlighting}
\end{Shaded}

Le nom de la fonction est \texttt{print}. L\textquotesingle expression
entre parenthèses est appelée l\textquotesingle argument de la fonction.
Le résultat, pour cette fonction, est le type de
l\textquotesingle argument.

Il est courant de dire qu\textquotesingle une fonction « prend » un
argument et « retourne » un résultat. Le résultat est appelé la valeur
de retour (\emph{return value}).

\hypertarget{fonctions-de-conversion-de-type}{%
\subsection{Fonctions de conversion de
type}\label{fonctions-de-conversion-de-type}}

Python fournit des fonctions intégrées qui convertissent les valeurs
d\textquotesingle un type à un autre. La fonction \texttt{int} prend
n\textquotesingle importe quelle valeur et la convertit en entier, si
elle le peut, ou signale une erreur sinon :

\begin{Shaded}
\begin{Highlighting}[]
\OperatorTok{\textgreater{}\textgreater{}\textgreater{}} \BuiltInTok{int}\NormalTok{(}\StringTok{\textquotesingle{}32\textquotesingle{}}\NormalTok{)}
\DecValTok{32}
\OperatorTok{\textgreater{}\textgreater{}\textgreater{}} \BuiltInTok{int}\NormalTok{(}\StringTok{\textquotesingle{}Hello\textquotesingle{}}\NormalTok{)}
\NormalTok{Traceback (most recent call last):}
\NormalTok{  File }\StringTok{"\textless{}stdin\textgreater{}"}\NormalTok{, line }\DecValTok{1}\NormalTok{, }\KeywordTok{in} \OperatorTok{\textless{}}\NormalTok{module}\OperatorTok{\textgreater{}}
\PreprocessorTok{ValueError}\NormalTok{: invalid literal }\ControlFlowTok{for} \BuiltInTok{int}\NormalTok{() }\ControlFlowTok{with}\NormalTok{ base }\DecValTok{10}\NormalTok{: }\StringTok{\textquotesingle{}Hello\textquotesingle{}}
\end{Highlighting}
\end{Shaded}

\texttt{int} peut convertir des valeurs à virgule flottante en entiers,
mais elle n\textquotesingle arrondit pas ; elle coupe simplement la
partie fractionnaire :

\begin{Shaded}
\begin{Highlighting}[]
\OperatorTok{\textgreater{}\textgreater{}\textgreater{}} \BuiltInTok{int}\NormalTok{(}\FloatTok{3.99999}\NormalTok{)}
\DecValTok{3}
\OperatorTok{\textgreater{}\textgreater{}\textgreater{}} \BuiltInTok{int}\NormalTok{(}\OperatorTok{{-}}\FloatTok{2.3}\NormalTok{)}
\OperatorTok{{-}}\DecValTok{2}
\end{Highlighting}
\end{Shaded}

La fonction \texttt{float} convertit les entiers et les chaînes de
caractères en nombres à virgule flottante :

\begin{Shaded}
\begin{Highlighting}[]
\OperatorTok{\textgreater{}\textgreater{}\textgreater{}} \BuiltInTok{float}\NormalTok{(}\DecValTok{32}\NormalTok{)}
\FloatTok{32.0}
\OperatorTok{\textgreater{}\textgreater{}\textgreater{}} \BuiltInTok{float}\NormalTok{(}\StringTok{\textquotesingle{}3.14159\textquotesingle{}}\NormalTok{)}
\FloatTok{3.14159}
\end{Highlighting}
\end{Shaded}

Enfin, \texttt{str} convertit son argument en une chaîne de caractères :

\begin{Shaded}
\begin{Highlighting}[]
\OperatorTok{\textgreater{}\textgreater{}\textgreater{}} \BuiltInTok{str}\NormalTok{(}\DecValTok{32}\NormalTok{)}
\CommentTok{\textquotesingle{}32\textquotesingle{}}
\OperatorTok{\textgreater{}\textgreater{}\textgreater{}} \BuiltInTok{str}\NormalTok{(}\FloatTok{3.14159}\NormalTok{)}
\CommentTok{\textquotesingle{}3.14159\textquotesingle{}}
\end{Highlighting}
\end{Shaded}

\hypertarget{fonctions-mathuxe9matiques}{%
\subsection{Fonctions mathématiques}\label{fonctions-mathuxe9matiques}}

Python possède un module \texttt{math} qui fournit la plupart des
fonctions mathématiques familières. Un module est un fichier qui
contient une collection de fonctions liées.

Avant de pouvoir utiliser les fonctions d\textquotesingle un module,
nous devons l\textquotesingle importer avec une instruction
\texttt{import} :

\begin{Shaded}
\begin{Highlighting}[]
\OperatorTok{\textgreater{}\textgreater{}\textgreater{}} \ImportTok{import}\NormalTok{ math}
\end{Highlighting}
\end{Shaded}

Cette instruction crée un objet module nommé \texttt{math}. Pour accéder
à l\textquotesingle une des fonctions, vous devez spécifier le nom du
module et le nom de la fonction, séparés par un point. Ce format est
appelé la notation pointée (\emph{dot notation}).

\begin{Shaded}
\begin{Highlighting}[]
\OperatorTok{\textgreater{}\textgreater{}\textgreater{}}\NormalTok{ ratio }\OperatorTok{=}\NormalTok{ puissance\_signal }\OperatorTok{/}\NormalTok{ puissance\_bruit}
\OperatorTok{\textgreater{}\textgreater{}\textgreater{}}\NormalTok{ decibels }\OperatorTok{=} \DecValTok{10} \OperatorTok{*}\NormalTok{ math.log10(ratio)}
\OperatorTok{\textgreater{}\textgreater{}\textgreater{}}\NormalTok{ radians }\OperatorTok{=} \FloatTok{0.7}
\OperatorTok{\textgreater{}\textgreater{}\textgreater{}}\NormalTok{ hauteur }\OperatorTok{=}\NormalTok{ math.sin(radians)}
\end{Highlighting}
\end{Shaded}

Le premier exemple utilise \texttt{math.log10} pour calculer un rapport
signal/bruit en décibels (en supposant que \texttt{puissance\_signal} et
\texttt{puissance\_bruit} sont définies). Le deuxième exemple trouve le
sinus de \texttt{radians}. Le nom de la variable est un indice indiquant
que \texttt{sin} et les autres fonctions trigonométriques (\texttt{cos},
\texttt{tan}, etc.) prennent des arguments en radians.

Pour convertir des degrés en radians, divisez par 180 et multipliez par
π :

\begin{Shaded}
\begin{Highlighting}[]
\OperatorTok{\textgreater{}\textgreater{}\textgreater{}}\NormalTok{ degres }\OperatorTok{=} \DecValTok{45}
\OperatorTok{\textgreater{}\textgreater{}\textgreater{}}\NormalTok{ radians }\OperatorTok{=}\NormalTok{ degres }\OperatorTok{/} \FloatTok{180.0} \OperatorTok{*}\NormalTok{ math.pi}
\OperatorTok{\textgreater{}\textgreater{}\textgreater{}}\NormalTok{ math.sin(radians)}
\FloatTok{0.7071067811865476}
\end{Highlighting}
\end{Shaded}

L\textquotesingle expression \texttt{math.pi} récupère la variable
\texttt{pi} depuis le module \texttt{math}. Sa valeur est une
approximation de π à virgule flottante, précise à environ 15 chiffres.

\hypertarget{composition}{%
\subsection{Composition}\label{composition}}

Jusqu\textquotesingle à présent, nous avons examiné les éléments
d\textquotesingle un programme --- variables, expressions et
instructions --- de manière isolée, sans parler de la façon de les
combiner.

L\textquotesingle une des fonctionnalités les plus utiles des langages
de programmation est leur capacité à prendre de petits blocs de
construction et à les composer. Par exemple, l\textquotesingle argument
d\textquotesingle une fonction peut être n\textquotesingle importe quel
type d\textquotesingle expression, y compris des opérateurs
arithmétiques :

\begin{Shaded}
\begin{Highlighting}[]
\NormalTok{x }\OperatorTok{=}\NormalTok{ math.sin(degres }\OperatorTok{/} \FloatTok{360.0} \OperatorTok{*} \DecValTok{2} \OperatorTok{*}\NormalTok{ math.pi)}
\end{Highlighting}
\end{Shaded}

Et même des appels de fonction :

\begin{Shaded}
\begin{Highlighting}[]
\NormalTok{x }\OperatorTok{=}\NormalTok{ math.exp(math.log(x}\OperatorTok{+}\DecValTok{1}\NormalTok{))}
\end{Highlighting}
\end{Shaded}

Presque partout où vous pouvez mettre une valeur, vous pouvez mettre une
expression arbitraire.

\hypertarget{ajout-de-nouvelles-fonctions}{%
\subsection{Ajout de nouvelles
fonctions}\label{ajout-de-nouvelles-fonctions}}

Jusqu\textquotesingle à présent, nous n\textquotesingle avons utilisé
que les fonctions intégrées à Python, mais il est également possible
d\textquotesingle ajouter de nouvelles fonctions. Une définition de
fonction spécifie le nom d\textquotesingle une nouvelle fonction et la
séquence d\textquotesingle instructions qui s\textquotesingle exécute
lorsque la fonction est appelée.

Voici un exemple :

\begin{Shaded}
\begin{Highlighting}[]
\KeywordTok{def}\NormalTok{ afficher\_paroles():}
    \BuiltInTok{print}\NormalTok{(}\StringTok{"Je suis un bûcheron, et je vais bien."}\NormalTok{)}
    \BuiltInTok{print}\NormalTok{(}\StringTok{"Je dors toute la nuit et je travaille toute la journée."}\NormalTok{)}
\end{Highlighting}
\end{Shaded}

\texttt{def} est un mot-clé qui indique qu\textquotesingle il
s\textquotesingle agit d\textquotesingle une définition de fonction. Le
nom de la fonction est \texttt{afficher\_paroles}. Les règles pour les
noms de fonctions sont les mêmes que pour les variables : les lettres,
les chiffres et le soulignement sont légaux, mais le premier caractère
ne peut pas être un chiffre.

Les parenthèses vides après le nom indiquent que cette fonction ne prend
aucun argument.

La première ligne de la définition de la fonction est appelée
l\textquotesingle en-tête (\emph{header}) ; le reste est appelé le corps
(\emph{body}). Le corps doit être indenté. Par convention,
l\textquotesingle indentation est toujours de quatre espaces.

La syntaxe pour appeler la nouvelle fonction est la même que pour les
fonctions intégrées :

\begin{Shaded}
\begin{Highlighting}[]
\OperatorTok{\textgreater{}\textgreater{}\textgreater{}}\NormalTok{ afficher\_paroles()}
\NormalTok{Je suis un bûcheron, et je vais bien.}
\NormalTok{Je dors toute la nuit et je travaille toute la journée.}
\end{Highlighting}
\end{Shaded}

Une fois que vous avez défini une fonction, vous pouvez
l\textquotesingle utiliser à l\textquotesingle intérieur
d\textquotesingle une autre fonction.

\begin{Shaded}
\begin{Highlighting}[]
\KeywordTok{def}\NormalTok{ repeter\_paroles():}
\NormalTok{    afficher\_paroles()}
\NormalTok{    afficher\_paroles()}
\end{Highlighting}
\end{Shaded}

\hypertarget{duxe9finitions-et-utilisations}{%
\subsection{Définitions et
utilisations}\label{duxe9finitions-et-utilisations}}

Rassemblez les fragments de code de la section précédente et
l\textquotesingle ensemble du programme ressemblera à ceci :

\begin{Shaded}
\begin{Highlighting}[]
\KeywordTok{def}\NormalTok{ afficher\_paroles():}
    \BuiltInTok{print}\NormalTok{(}\StringTok{"Je suis un bûcheron, et je vais bien."}\NormalTok{)}
    \BuiltInTok{print}\NormalTok{(}\StringTok{"Je dors toute la nuit et je travaille toute la journée."}\NormalTok{)}

\KeywordTok{def}\NormalTok{ repeter\_paroles():}
\NormalTok{    afficher\_paroles()}
\NormalTok{    afficher\_paroles()}

\NormalTok{repeter\_paroles()}
\end{Highlighting}
\end{Shaded}

Ce programme contient deux définitions de fonctions :
\texttt{afficher\_paroles} et \texttt{repeter\_paroles}. Les définitions
de fonctions sont exécutées comme d\textquotesingle autres instructions,
mais l\textquotesingle effet est de créer des objets fonctions. Les
instructions à l\textquotesingle intérieur de la fonction ne sont pas
exécutées tant que la fonction n\textquotesingle est pas appelée.

Comme vous pouvez vous y attendre, vous devez créer une fonction avant
de pouvoir l\textquotesingle exécuter. Autrement dit, la définition de
la fonction doit être exécutée avant la première fois
qu\textquotesingle elle est appelée.

\hypertarget{flux-dexuxe9cution}{%
\subsection{Flux d\textquotesingle exécution}\label{flux-dexuxe9cution}}

Pour s\textquotesingle assurer qu\textquotesingle une fonction est
définie avant sa première utilisation, il faut connaître
l\textquotesingle ordre dans lequel les instructions sont exécutées, ce
qu\textquotesingle on appelle le flux d\textquotesingle exécution
(\emph{flow of execution}).

L\textquotesingle exécution commence toujours par la première
instruction du programme. Les instructions sont exécutées une à la fois,
de haut en bas.

Les définitions de fonctions ne modifient pas le flux
d\textquotesingle exécution, mais rappelez-vous que les instructions à
l\textquotesingle intérieur d\textquotesingle une fonction ne sont pas
exécutées tant que la fonction n\textquotesingle est pas appelée.

\hypertarget{paramuxe8tres-et-arguments}{%
\subsection{Paramètres et arguments}\label{paramuxe8tres-et-arguments}}

Certaines des fonctions intégrées que nous avons vues nécessitent des
arguments. Par exemple, lorsque vous appelez \texttt{math.sin}, vous
passez un nombre comme argument. Certaines fonctions prennent plus
d\textquotesingle un argument : \texttt{math.pow} prend deux arguments,
la base et l\textquotesingle exposant.

À l\textquotesingle intérieur de la fonction, les arguments sont
affectés à des variables appelées paramètres. Voici un exemple
d\textquotesingle une fonction définie par l\textquotesingle utilisateur
qui prend un argument :

\begin{Shaded}
\begin{Highlighting}[]
\KeywordTok{def}\NormalTok{ imprimer\_deux\_fois(bruce):}
    \BuiltInTok{print}\NormalTok{(bruce)}
    \BuiltInTok{print}\NormalTok{(bruce)}
\end{Highlighting}
\end{Shaded}

Cette fonction affecte l\textquotesingle argument à un paramètre nommé
\texttt{bruce}. Lorsque la fonction est appelée, elle imprime la valeur
du paramètre (quel qu\textquotesingle il soit) à deux reprises.

\begin{Shaded}
\begin{Highlighting}[]
\OperatorTok{\textgreater{}\textgreater{}\textgreater{}}\NormalTok{ imprimer\_deux\_fois(}\StringTok{\textquotesingle{}Spam\textquotesingle{}}\NormalTok{)}
\NormalTok{Spam}
\NormalTok{Spam}
\OperatorTok{\textgreater{}\textgreater{}\textgreater{}}\NormalTok{ imprimer\_deux\_fois(}\DecValTok{17}\NormalTok{)}
\DecValTok{17}
\DecValTok{17}
\OperatorTok{\textgreater{}\textgreater{}\textgreater{}}\NormalTok{ imprimer\_deux\_fois(math.pi)}
\FloatTok{3.141592653589793}
\FloatTok{3.141592653589793}
\end{Highlighting}
\end{Shaded}

\hypertarget{les-variables-et-paramuxe8tres-sont-locaux}{%
\subsection{Les variables et paramètres sont
locaux}\label{les-variables-et-paramuxe8tres-sont-locaux}}

Lorsque vous créez une variable à l\textquotesingle intérieur
d\textquotesingle une fonction, elle est locale, ce qui signifie
qu\textquotesingle elle n\textquotesingle existe qu\textquotesingle à
l\textquotesingle intérieur de cette fonction. Par exemple :

\begin{Shaded}
\begin{Highlighting}[]
\KeywordTok{def}\NormalTok{ concatener\_deux\_fois(partie1, partie2):}
\NormalTok{    chat }\OperatorTok{=}\NormalTok{ partie1 }\OperatorTok{+}\NormalTok{ partie2}
\NormalTok{    imprimer\_deux\_fois(chat)}
\end{Highlighting}
\end{Shaded}

Cette fonction prend deux arguments, les concatène et affiche le
résultat deux fois.

\begin{Shaded}
\begin{Highlighting}[]
\OperatorTok{\textgreater{}\textgreater{}\textgreater{}}\NormalTok{ ligne1 }\OperatorTok{=} \StringTok{\textquotesingle{}Bing tiddle \textquotesingle{}}
\OperatorTok{\textgreater{}\textgreater{}\textgreater{}}\NormalTok{ ligne2 }\OperatorTok{=} \StringTok{\textquotesingle{}tiddle bang.\textquotesingle{}}
\OperatorTok{\textgreater{}\textgreater{}\textgreater{}}\NormalTok{ concatener\_deux\_fois(ligne1, ligne2)}
\NormalTok{Bing tiddle tiddle bang.}
\NormalTok{Bing tiddle tiddle bang.}
\end{Highlighting}
\end{Shaded}

Lorsque la fonction se termine, la variable \texttt{chat} est détruite.
Si nous essayons de l\textquotesingle imprimer depuis
l\textquotesingle extérieur de la fonction, nous obtenons une exception
:

\begin{Shaded}
\begin{Highlighting}[]
\OperatorTok{\textgreater{}\textgreater{}\textgreater{}} \BuiltInTok{print}\NormalTok{(chat)}
\PreprocessorTok{NameError}\NormalTok{: name }\StringTok{\textquotesingle{}chat\textquotesingle{}} \KeywordTok{is} \KeywordTok{not}\NormalTok{ defined}
\end{Highlighting}
\end{Shaded}

Les paramètres sont également locaux. À l\textquotesingle extérieur de
\texttt{imprimer\_deux\_fois}, il n\textquotesingle y a pas de variable
telle que \texttt{bruce}.

\hypertarget{fonctions-productives-et-fonctions-uxe0-effet-de-bord}{%
\subsection{Fonctions productives et fonctions à effet de
bord}\label{fonctions-productives-et-fonctions-uxe0-effet-de-bord}}

Certaines des fonctions que nous utilisons, telles que les fonctions
mathématiques, retournent des résultats ; faute d\textquotesingle un
meilleur terme, je les appelle des fonctions productives (\emph{fruitful
functions}). D\textquotesingle autres fonctions, comme
\texttt{imprimer\_deux\_fois}, effectuent une action mais ne retournent
pas de valeur. Elles sont appelées fonctions nulles ou fonctions à effet
de bord (\emph{void functions}).

Si vous appelez une fonction productive en mode script ou dans un
programme, la valeur de retour est généralement ignorée, à moins que
vous ne l\textquotesingle affectiez à une variable :

\begin{Shaded}
\begin{Highlighting}[]
\NormalTok{math.cos(radians)}
\NormalTok{x }\OperatorTok{=}\NormalTok{ math.cos(radians)}
\end{Highlighting}
\end{Shaded}

Les fonctions \emph{void} peuvent afficher quelque chose à
l\textquotesingle écran ou avoir un autre effet, mais leur valeur de
retour est une valeur spéciale appelée \texttt{None}.

\begin{Shaded}
\begin{Highlighting}[]
\OperatorTok{\textgreater{}\textgreater{}\textgreater{}}\NormalTok{ resultat }\OperatorTok{=}\NormalTok{ imprimer\_deux\_fois(}\StringTok{\textquotesingle{}Bing\textquotesingle{}}\NormalTok{)}
\NormalTok{Bing}
\NormalTok{Bing}
\OperatorTok{\textgreater{}\textgreater{}\textgreater{}} \BuiltInTok{print}\NormalTok{(resultat)}
\VariableTok{None}
\end{Highlighting}
\end{Shaded}

La valeur \texttt{None} n\textquotesingle est pas la même chose que la
chaîne de caractères \texttt{\textquotesingle{}None\textquotesingle{}}.
C\textquotesingle est une valeur spéciale qui a son propre type :

\begin{Shaded}
\begin{Highlighting}[]
\OperatorTok{\textgreater{}\textgreater{}\textgreater{}} \BuiltInTok{type}\NormalTok{(}\VariableTok{None}\NormalTok{)}
\OperatorTok{\textless{}}\KeywordTok{class} \StringTok{\textquotesingle{}NoneType\textquotesingle{}}\OperatorTok{\textgreater{}}
\end{Highlighting}
\end{Shaded}

\hypertarget{pourquoi-des-fonctions}{%
\subsection{Pourquoi des fonctions ?}\label{pourquoi-des-fonctions}}

Il n\textquotesingle est pas toujours clair au premier abord pourquoi il
vaut la peine de diviser un programme en fonctions. Voici quelques
raisons :

\begin{itemize}
\tightlist
\item
  La création d\textquotesingle une nouvelle fonction vous donne
  l\textquotesingle opportunité de nommer un groupe
  d\textquotesingle instructions, ce qui rend votre programme plus
  facile à lire et à déboguer.
\item
  Les fonctions peuvent rendre un programme plus petit en éliminant le
  code répétitif. Plus tard, si vous apportez un changement, vous
  n\textquotesingle aurez à le faire qu\textquotesingle à un seul
  endroit.
\item
  La division d\textquotesingle un long programme en fonctions vous
  permet de déboguer les parties une à la fois, puis de les assembler
  dans un tout fonctionnel.
\item
  Les fonctions bien conçues sont souvent utiles pour de nombreux
  programmes. Une fois que vous en avez écrit et débogué une, vous
  pouvez la réutiliser.
\end{itemize}

\hypertarget{duxe9bogage}{%
\subsection{Débogage}\label{duxe9bogage}}

L\textquotesingle une des compétences les plus importantes que vous
allez acquérir est le débogage (\emph{debugging}). Si votre programme
comporte une erreur d\textquotesingle exécution, Python imprimera un
message d\textquotesingle erreur (\emph{traceback}), qui contient des
informations utiles. Il indique où se trouve l\textquotesingle erreur,
le type de l\textquotesingle erreur, et quelles fonctions étaient en
cours d\textquotesingle exécution au moment de l\textquotesingle erreur.

\hypertarget{glossaire}{%
\subsection{Glossaire}\label{glossaire}}

fonction function Une séquence nommée d\textquotesingle instructions qui
effectue une opération utile. Les fonctions peuvent ou non prendre des
arguments, et peuvent ou non produire un résultat.

définition de fonction function definition Une instruction qui crée une
nouvelle fonction, en spécifiant son nom, ses paramètres et les
instructions qu\textquotesingle elle exécute.

en-tête header La première ligne d\textquotesingle une définition de
fonction.

corps body La séquence d\textquotesingle instructions à
l\textquotesingle intérieur d\textquotesingle une définition de
fonction.

paramètre parameter Un nom utilisé à l\textquotesingle intérieur
d\textquotesingle une fonction pour faire référence à la valeur passée
comme argument.

appel de fonction function call Une instruction qui exécute une
fonction. Elle se compose du nom de la fonction suivi
d\textquotesingle une liste d\textquotesingle arguments entre
parenthèses.

\begin{description}
\item[argument]
Une valeur fournie à une fonction lorsque la fonction est appelée. Cette
valeur est assignée au paramètre correspondant.
\end{description}

variable locale local variable Une variable définie à
l\textquotesingle intérieur d\textquotesingle une fonction. Une variable
locale ne peut être utilisée qu\textquotesingle à
l\textquotesingle intérieur de sa fonction.

valeur de retour return value Le résultat d\textquotesingle une
fonction.

fonction productive fruitful function Une fonction qui retourne une
valeur.

fonction nulle void function Une fonction qui ne retourne aucune valeur
(elle retourne toujours \texttt{None}).

\begin{description}
\item[module]
Un fichier qui contient une collection de fonctions et
d\textquotesingle autres définitions.
\end{description}

notation pointée dot notation La syntaxe pour appeler une fonction dans
un autre module, en spécifiant le nom du module suivi
d\textquotesingle un point et du nom de la fonction.

flux d\textquotesingle exécution flow of execution
L\textquotesingle ordre dans lequel les instructions sont exécutées lors
du fonctionnement d\textquotesingle un programme.

\hypertarget{exercices}{%
\subsection{Exercices}\label{exercices}}

\textbf{Exercice 1}

Écrivez une fonction nommée \texttt{right\_justify} qui prend une chaîne
de caractères nommée \texttt{s} comme paramètre et imprime la chaîne
avec suffisamment d\textquotesingle espaces de début pour que la
dernière lettre de la chaîne se trouve dans la colonne 70 de
l\textquotesingle affichage.

\begin{Shaded}
\begin{Highlighting}[]
\OperatorTok{\textgreater{}\textgreater{}\textgreater{}}\NormalTok{ right\_justify(}\StringTok{\textquotesingle{}monty\textquotesingle{}}\NormalTok{)}
\NormalTok{                                                                 monty}
\end{Highlighting}
\end{Shaded}

\emph{Indice : Utilisez la concaténation de chaînes et la répétition. De
plus, Python fournit une fonction intégrée appelée len qui renvoie la
longueur d\textquotesingle une chaîne, donc la valeur de
len(\textquotesingle monty\textquotesingle) est 5.}

\textbf{Exercice 2}

Un objet de fonction est une valeur que vous pouvez affecter à une
variable ou passer comme argument. Par exemple, \texttt{do\_twice} est
une fonction qui prend un objet de fonction comme argument et
l\textquotesingle appelle à deux reprises :

\begin{Shaded}
\begin{Highlighting}[]
\KeywordTok{def}\NormalTok{ do\_twice(f):}
\NormalTok{    f()}
\NormalTok{    f()}
\end{Highlighting}
\end{Shaded}

Voici un exemple qui utilise \texttt{do\_twice} pour appeler une
fonction nommée \texttt{print\_spam} à deux reprises.

\begin{Shaded}
\begin{Highlighting}[]
\KeywordTok{def}\NormalTok{ print\_spam():}
    \BuiltInTok{print}\NormalTok{(}\StringTok{\textquotesingle{}spam\textquotesingle{}}\NormalTok{)}

\NormalTok{do\_twice(print\_spam)}
\end{Highlighting}
\end{Shaded}

\begin{enumerate}
\def\labelenumi{\arabic{enumi}.}
\tightlist
\item
  Tapez cet exemple dans un script et testez-le.
\item
  Modifiez \texttt{do\_twice} afin qu\textquotesingle elle prenne deux
  arguments, un objet fonction et une valeur, et appelle la fonction
  deux fois, en passant la valeur comme argument.
\end{enumerate}
